% Who Should Review Whom? Response-Conditioned Disassortative Peer Review
% Compile: pdflatex main.tex (requires a LaTeX distribution)
\documentclass[11pt]{article}
\usepackage[margin=1in]{geometry}
\usepackage{amsmath,amssymb}
\usepackage{booktabs}
\usepackage{graphicx}
\usepackage{hyperref}
\usepackage{algorithm}
\usepackage{algpseudocode}
\usepackage{caption}

\newcommand{\TBD}[1]{\textbf{[TBD: #1]}}

\title{Who Should Review Whom?\\Response-Conditioned Disassortative Peer Review for Test-Time LLM Reasoning}
\author{Anonymous}
\date{September 2026}

\begin{document}
\maketitle

\begin{abstract}
\TBD{filled after full-study results are frozen. Problem: homogeneous same-prompt LLM samples have correlated errors; method: Farthest-Pair Reciprocal Review (FPRR), a training-free response-conditioned review topology; setup: controlled comparison against random/nearest pairing, maximum-weight matching, self-review, and majority voting on identical cached initial responses; strongest result: \TBD{}; limitation: \TBD{}.}
\end{abstract}

\section{Introduction}
\label{sec:intro}

Repeated stochastic inference from a fixed language model can substantially improve answer accuracy without touching the model's weights: self-consistency voting over independent samples~\cite{wang2022selfconsistency} and sample-and-vote scaling~\cite{li2024moreagents} are the canonical examples. A natural next step is to let sampled instances \emph{communicate}: multi-agent debate~\cite{du2023multiagent} and peer-review schemes~\cite{xu2023peerreview} allow instances of a model to critique each other's answers before aggregation.

Yet communication among homogeneous agents has a structural weakness. Instances of the same checkpoint, given the same prompt, do not err independently: they share training data, inductive biases, and failure modes, so their samples form \emph{correlated reasoning clusters} — groups of answers that share assumptions, strategies, and misconceptions. Broadcast-style communication then tends to reinforce the majority view, which is exactly the view most likely to embody the shared error. Empirically, debate often fails to beat simple voting~\cite{yang2025revisiting,choi2025debateorvote}, and diversity-aware variants that prune or retain messages still communicate by broadcast~\cite{estornell2024multillm,nguyen2026dars}.

This paper asks a question that, to our knowledge, has not been isolated experimentally:

\begin{quote}
\emph{Who should review whom?} Given $N$ independent answers from the same model, does the \textbf{topology} of pairwise peer review — specifically, pairing each agent with its most semantically \emph{dissimilar} peer — extract more correct answers than random pairing, similarity-based pairing, debate, or voting, under identical model, prompt, and inference budget?
\end{quote}

We instantiate this question as \textbf{Farthest-Pair Reciprocal Review (FPRR)}: embed each response (conclusion + concise justification), compute a pairwise cosine-distance matrix, construct a perfect matching that pairs distant responses (a randomized greedy farthest matching, RGFM, with an exact maximum-weight matching as a variant), run reciprocal pair-limited peer review within each pair, and synthesize a final answer from all answers and reviews. The pipeline is training-free, retrieval-free, and uses one frozen model and one frozen prompt set throughout.

Crucially, our design isolates the causal contribution of the routing rule: all pairing conditions consume the \emph{exact same cached initial responses}, so differences cannot be attributed to resampling luck; and a compute-matched self-review control separates topology effects from merely spending more inference.

\paragraph{Contributions.}
\begin{itemize}
\item We formalize response-conditioned review topology and provide the first controlled study of disassortative (farthest-pair) reciprocal review, including exact maximum-weight matching, randomized greedy matching, nearest pairing, and random pairing on identical response sets.
\item \TBD{main empirical claim, after results}
\item A mechanistic analysis of what information crosses disassortative edges: pair-composition effects, distance-vs-correction curves, majority-error recovery, correct-majority corruption, and quantified failure modes (outlier amplification, shared misconception, judge failure).
\end{itemize}

\section{Related Work}
\label{sec:related}

\paragraph{Test-time scaling and sampling ensembles.}
Self-consistency~\cite{wang2022selfconsistency} aggregates diverse chains of thought by majority vote; Li et al.~\cite{li2024moreagents} show accuracy scales with agent count under voting alone (their similarity measure selects the \emph{most} similar sample — the opposite polarity to ours). Embedding-similarity pruning of single-agent trace pools (Slim-SC~\cite{slimsc}, SSDP~\cite{ssdp}, DeepPrune~\cite{deepprune}) improves efficiency but constructs no inter-agent communication. DIVSE~\cite{divse} diversifies at the prompt level. None condition a \emph{communication topology} on response-level distances.

\paragraph{Multi-agent debate.}
Du et al.~\cite{du2023multiagent} established broadcast debate among homogeneous instances; ReConcile~\cite{chen2023reconcile} uses heterogeneous models with confidence-weighted consensus. Controlled replications temper the claims: debate's benefit is conditional and often matches self-agent baselines~\cite{yang2025revisiting}; voting can dominate debate~\cite{choi2025debateorvote}; and agents' debate behavior itself is questionable under controlled logical tasks~\cite{wu2025reallydebate}. All use broadcast or fixed topologies.

\paragraph{Diversity in debate.}
Estornell \& Liu~\cite{estornell2024multillm} compute response embeddings and pairwise distances, diagnose echo chambers, and prune to a max-diversity subset — but survivors still broadcast; there is no pairing or reciprocal review. DAR~\cite{nguyen2026dars} retains a disagreeing subset for broadcast via an LLM filter (embeddings appear only as an evaluation metric). Demystifying MAD~\cite{zhu2026demystifying} initializes pools by maximizing unique-answer \emph{counts} and trains confidence via GRPO. S$^2$-MAD~\cite{zeng2025s2mad} uses response cosine similarity with the opposite polarity (suppressing redundant similar viewpoints). RMoA~\cite{rmoa2025} performs farthest-point-style selection of responses as \emph{aggregation inputs}, not review partners. Diversity in this literature prunes, weights, retains, or initializes — it never builds a one-to-one review graph.

\paragraph{Distance-conditioned topology.}
CortexDebate~\cite{cortexdebate2025}, the nearest work to ours, includes output cosine dissimilarity in sparse debate-graph edge weights (trust-formula composite, heterogeneous models, thresholded broadcast graph, majority vote). DySCo~\cite{dysco2026} lists reasoning-embedding cosine distance as one of five optional edge signals with one-way top-$k$ critique. DyTopo~\cite{dytopo2026} routes by \emph{similarity} of need/offer descriptors. MACE~\cite{mace2026} conditions peer selection on n-gram diversity via a bandit. Learned-topology methods (GPTSwarm~\cite{gptswarm}, DyLAN~\cite{dylan}, AgentPrune~\cite{agentprune}, G-Designer~\cite{gdesigner}, MacNet~\cite{macnet}, Graph-GRPO~\cite{graphgrpo}, Guided Topology Diffusion~\cite{gtd2025}, Codebook Agent~\cite{codebookagent}) condition on queries, roles, or learned values — never on distances among current responses — and usually require training.

\paragraph{Peer review and pairing.}
Xu et al.~\cite{xu2023peerreview} implement reciprocal review on an all-to-all graph; PiCO~\cite{pico2024} pairs agents for mutual assessment but pairs them at random (in effect, our random-pairing control); a tournament scheme~\cite{qdtournament} pairs by skill standing. A 141-study survey~\cite{madsurvey2026} reports 88.7\% static topologies and no embedding-distance pairing anywhere in its taxonomy. The empty intersection our method occupies — response-distance $\rightarrow$ farthest perfect matching $\rightarrow$ reciprocal pair-limited review $\rightarrow$ synthesis, all training-free over homogeneous same-prompt samples — is documented in our novelty-gate report (\texttt{NOVELTY\_GATE.md}), which examined 40+ papers with full-text method-section reading and citation snowballing over 2{,}235 citing papers of~\cite{du2023multiagent}.

\section{Method}
\label{sec:method}

\subsection{Problem setting}
For each problem $x$, $N$ (even) homogeneous agents — the same checkpoint, system prompt, problem text, and decoding configuration — independently produce responses
$r_i = (\text{answer}_i, \text{justification}_i)$, $i=1,\dots,N$,
where the justification is a concise visible reasoning summary (no hidden chain-of-thought is requested or used). Independent stochastic sampling is the only source of diversity.

\subsection{Response representation and distance}
Each response is embedded as $e_i = f(\text{CONCLUSION: answer}_i \mathbin{\Vert} \text{JUSTIFICATION: justification}_i)$ with a pinned sentence encoder $f$ (bge-m3, 1024-d, Ollama digest \texttt{790764642607}). The question text is never embedded, since identical question text would dominate similarity. The primary distance is cosine distance $d_{ij} = 1 - \cos(e_i, e_j)$ (metric D1). Ablations: D2 (justification-only embeddings), D3 (hybrid $\lambda d_{ij} + (1-\lambda)\mathbf{1}[a_i \neq a_j]$, $\lambda$ tuned on pilot only), D4 (symmetrized NLI contradiction probability, roberta-large-mnli).

\subsection{Pairing algorithms}
\begin{algorithm}[t]
\caption{Randomized Greedy Farthest Matching (RGFM)}
\label{alg:rgfm}
\begin{algorithmic}[1]
\Require distance matrix $d \in \mathbb{R}^{N\times N}$, seeded PRNG
\State $U \gets \{1,\dots,N\}$; $M \gets \emptyset$
\While{$U \neq \emptyset$}
  \State draw $i$ uniformly from $U$ (seeded)
  \State $j \gets \arg\max_{k \in U \setminus \{i\}} d_{ik}$
  \State $M \gets M \cup \{(i,j)\}$; $U \gets U \setminus \{i,j\}$
\EndWhile
\State \Return perfect matching $M$
\end{algorithmic}
\end{algorithm}

FPRR uses RGFM (Algorithm~\ref{alg:rgfm}); the seed is recorded because the greedy rule is order-dependent. Controls: \textbf{Random} (seeded random perfect matching), \textbf{Nearest} (same loop with $\arg\min$), and \textbf{MaxWeight} (exact maximum-weight perfect matching $M^* = \arg\max_M \sum_{(i,j)\in M} d_{ij}$ by exhaustive enumeration over the $(N-1)!!$ matchings; $N \le 8$).

\subsection{Reciprocal peer review}
For every pair $(i,j) \in M$, agent $i$ reviews $j$ and $j$ reviews $i$. The reviewer sees the problem, its own answer and justification, and the partner's answer and justification — nothing else: no other agents' answers, no ground truth, no embedding distances, no indication of how the pair was formed. The review prompt (frozen, identical across pairing conditions) requires the reviewer to check the peer independently, list correct claims, specific errors, and unverified assumptions, and end with a verdict (\texttt{PEER CORRECT}/\texttt{PEER INCORRECT}) and an updated own answer. The wording explicitly forbids assuming either side is correct, to counter both sycophancy and reflexive disagreement.

\subsection{Final synthesis}
One synthesizer call (same model) receives the problem, all anonymized candidate answers (labels S1\dots SN after a seeded per-question shuffle), and — for review conditions — all pairwise reviews. It never sees ground truth, condition names, distances, or pairing information. Its instructions: do not count repeated claims as independent evidence; evaluate argument quality; use criticisms to locate errors; recover correct minority solutions when justified; ignore rhetorical confidence. The same synthesizer prompt is used in all review conditions.

\subsection{Compute accounting}
Per question, FPRR costs $N$ initial generations $+$ $N$ review calls $+$ $1$ synthesis call. The compute-matched control (C9) spends the identical budget on $N$ self-reviews $+$ $1$ synthesis over the same cached initials.

\section{Experimental Setup}
\label{sec:setup}

\paragraph{Model.} All agents and the synthesizer use one frozen served model: \texttt{deepseek-flash} (DeepSeek-V4.1-Flash) via the DeepSeek API, temperature $0.8$, top-$p$ $0.95$, max tokens 1500 (solve/review) / 2500 (synthesis), thinking mode disabled so that the visible justification is the only reasoning artifact; no seed parameter is exposed by the provider. Stochasticity was verified empirically (distinct outputs on repeated identical calls). The model is homogeneous across all agents and conditions; no weights are modified.

\paragraph{Benchmarks.} MMLU-Pro~\cite{mmlupro} (test, 12{,}032 items, CC-BY-4.0) and SuperGPQA~\cite{supergpqa} (26{,}526 items, ODC-BY; sampled) as primaries; GSM8K~\cite{gsm8k} (test) as a sanity check (pilot: 95\% single-agent $\Rightarrow$ ceiling-saturated, excluded from full runs); MATH-500 reserved. GPQA-Diamond and LiveBench were inaccessible from this environment (access-gated); SuperGPQA serves as the recent graduate-level benchmark. Question IDs are fixed by a stratified sampler (sampling seed 42) and shared across all conditions and seeds.

\paragraph{Conditions.} C0 single response; C1 majority vote over $N{=}4$ cached responses; C2 cached responses $\rightarrow$ synthesizer (no review); C3/C4/C5/C6 = reciprocal review under Random / Nearest / RGFM / MaxWeight pairings on the \emph{same} cached responses; C9 compute-matched self-review; C7 one-round broadcast debate (pilot only). All pairing conditions consume identical initial responses per (question, seed).

\paragraph{Sample sizes and seeds.} Pilot: 100 questions per benchmark. Full: MMLU-Pro 500 questions $\times$ 3 experimental seeds; SuperGPQA 500 $\times$ 2 seeds. Ablations on 200-question subsets: $N \in \{2,4,6,8\}$; pairing seeds $\{1,2\}$ for Random and RGFM; distance metrics D2/D3/D4.

\paragraph{Statistics.} Exact-answer accuracy with a frozen deterministic parser (last \texttt{FINAL ANSWER:} line; unparsed $\Rightarrow$ incorrect). Paired bootstrap (10{,}000 resamples over questions) for 95\% CIs of accuracy differences; continuity-corrected McNemar tests for C5 vs.\ \{C3, C4, C1, C9\} with Holm--Bonferroni correction. Hypotheses H1--H6 were pre-registered (see \texttt{preregistration.md}, including all amendments) before any full-test evaluation.

\paragraph{Contamination.} GSM8K and MATH are old and heavily exposed; we therefore base no claim on them. MMLU-Pro (June 2024) and especially SuperGPQA (January 2025) postdate much training data; SuperGPQA's graduate-level items and the model's mid-band accuracy there (61\% single-agent in pilot) make it the more sensitive instrument. We nevertheless treat residual contamination as a limitation: it can only inflate absolute numbers, and applies identically to every condition, so paired differences remain interpretable.

\paragraph{Reproducibility.} Every API call is logged to JSONL (question ID, full prompt, raw output, parsed fields, token usage, latency, timestamp, condition, seeds). Initial responses are cached and shared across conditions; pairing/matching decisions depend only on distances, never on gold labels. Code and configs are versioned (git); analysis scripts are deterministic (seeded bootstrap).

% ================= RESULTS (filled after full study) =================

\section{Main Results}
\label{sec:results}
\TBD{Table: per-benchmark per-condition accuracy (pooled + per seed), paired bootstrap CIs, McNemar/Holm for C5 vs C3/C4/C1/C9; accuracy-per-token table; verdict on H1--H3.}

\section{Ablation Studies}
\label{sec:ablations}
\TBD{Pairing rule (C3--C6 same initials); agent count N in \{2,4,6,8\}; distance metric D1--D4; review vs no review (C2); synthesizer vs vote (C1); compute matching (C9); pairing-seed variance.}

\section{Mechanistic Analysis}
\label{sec:mechanism}
\TBD{Distance-decile vs correction probability; pair composition; majority-error recovery; minority rescue; consensus transitions; what crossed the disassortative edges.}

\section{Case Studies}
\label{sec:cases}
\TBD{successful correction; minority rescue; harmful outlier pairing (or its documented absence); shared-misconception failure.}

\section{Discussion}
\label{sec:discussion}
\TBD{interpretation conditional on results; system-level inference gains are not intrinsic model improvements; cluster/echo-chamber model of homogeneous errors if supported.}

\section{Limitations}
\label{sec:limitations}
Single served model (one provider, one checkpoint); API nondeterminism not seed-pinnable; benchmark contamination cannot be excluded; embedding dependence (bge-m3); synthesizer/judge bias; compute overhead relative to voting; low answer diversity at temperature 0.8 limits the exploitable disagreement mass; two benchmarks only; cross-seed observations on shared questions are mildly dependent.

\section{Conclusion}
\label{sec:conclusion}
\TBD{Result category A--F per the pre-registered decision criteria.}

% ================= REFERENCES =================
\begin{thebibliography}{99}

\bibitem{wang2022selfconsistency} Wang, X. et al. Self-Consistency Improves Chain of Thought Reasoning in Language Models. ICLR 2023 (arXiv:2203.11171).
\bibitem{du2023multiagent} Du, Y. et al. Improving Factuality and Reasoning in Language Models through Multiagent Debate. ICML 2024 (arXiv:2305.14325).
\bibitem{chen2023reconcile} Chen, J. et al. ReConcile: Round-Table Conference Improves Reasoning via Consensus among Diverse LLMs. ACL 2024 (arXiv:2309.13007).
\bibitem{li2024moreagents} Li, J. et al. More Agents Is All You Need. arXiv:2402.05120, 2024.
\bibitem{estornell2024multillm} Estornell, A. and Liu, Y. Multi-LLM Debate: Framework, Principals, and Interventions. NeurIPS 2024 (OpenReview sy7eSEXdPC).
\bibitem{yang2025revisiting} Yang, Z. et al. Revisiting Multi-Agent Debate as Test-Time Scaling: A Systematic Study of Conditional Effectiveness. arXiv:2505.22960, 2025.
\bibitem{choi2025debateorvote} Choi, Zhu, Li. Debate or Vote: Which Yields Better Decisions in Multi-Agent Large Language Models? NeurIPS 2025 (arXiv:2508.17536).
\bibitem{wu2025reallydebate} Wu, et al. Can LLM Agents Really Debate? A Controlled Study of Multi-Agent Debate in Logical Reasoning. arXiv:2511.07784, 2025.
\bibitem{zhu2026demystifying} Zhu, et al. Demystifying Multi-Agent Debate: The Role of Confidence and Diversity. arXiv:2601.19921, 2026.
\bibitem{nguyen2026dars} Nguyen, et al. Hear Both Sides: Efficient Multi-Agent Debate via Diversity-Aware Message Retention. arXiv:2603.20640, 2026.
\bibitem{zeng2025s2mad} Zeng, et al. S$^2$-MAD: Breaking the Token Barrier to Enhance Multi-Agent Debate Efficiency. NAACL 2025 (arXiv:2502.04790).
\bibitem{rmoa2025} Xie, et al. RMoA: Optimizing Mixture-of-Agents through Diversity Maximization and Residual Compensation. Findings of ACL 2025.
\bibitem{cortexdebate2025} Sun, et al. CortexDebate: Debating with Sparse, Dynamic and Trust-Aware Topology. Findings of ACL 2025 (arXiv:2507.03928).
\bibitem{dysco2026} Gou and Liu. DySCo: Dynamic Trust-Aware Sparse Consensus. arXiv:2606.01828, 2026.
\bibitem{dytopo2026} DyTopo: Dynamic Topology Routing for Multi-Agent Systems. arXiv:2602.06039, 2026.
\bibitem{mace2026} Li, et al. Multi-Agent LLMs Fail to Explore Each Other / MACE. arXiv:2607.11250, 2026.
\bibitem{gptswarm} Zhuge, et al. GPTSwarm: Language Agents as Optimizable Graphs. ICML 2024 (arXiv:2402.16823).
\bibitem{dylan} Liu, et al. Dynamic LLM-Agent Network (DyLAN). COLM 2024 (arXiv:2310.02170).
\bibitem{agentprune} Zhang, et al. AgentPrune: Reducing Redundancy in Multi-Agent Communication. ICLR 2025 (arXiv:2410.02506).
\bibitem{gdesigner} Zhang, et al. G-Designer: Architecting Multi-agent Communication Topologies via Graph Neural Networks. arXiv:2410.11782, 2024.
\bibitem{macnet} Qian, et al. MacNet: Multi-Agent Collaboration Networks. arXiv:2406.07155, 2024.
\bibitem{graphgrpo} Graph-GRPO: Graph Topology Optimization via GRPO. arXiv:2603.02701, 2026.
\bibitem{gtd2025} Jiang, et al. Guided Topology Diffusion for Multi-Agent Systems. ACL 2026 (arXiv:2510.07799).
\bibitem{codebookagent} Codebook Agent: Query-Conditioned Topology Prediction. arXiv:2609.02264, 2026.
\bibitem{xu2023peerreview} Xu, et al. Towards Reasoning in LLMs via Multi-Agent Peer Review Collaboration. arXiv:2311.08152, 2023.
\bibitem{pico2024} PiCO: Peer-based Collaborative Optimization. arXiv:2402.01830, 2024.
\bibitem{qdtournament} Quality-Diversity Debate Tournaments. arXiv:2510.05909, 2025.
\bibitem{madsurvey2026} Motger, et al. A Survey of Multi-Agent Debate Systems (141 studies). arXiv:2607.26212, 2026.
\bibitem{slimsc} Slim-SC: Efficient Self-Consistency via Similarity Pruning. arXiv:2509.13990, 2025.
\bibitem{ssdp} SSDP: Semantic-Similarity-Based Diverse Pruning. arXiv:2511.08595, 2025.
\bibitem{deepprune} DeepPrune: Efficient Trace Pruning for Test-Time Scaling. arXiv:2510.08483, 2025.
\bibitem{divse} Naik, et al. DIVSE: Diverse Self-Consistency. arXiv:2310.07088, 2023.
\bibitem{spmad} Li, et al. Sparse Communication Topology in Multi-Agent Debate. Findings of EMNLP 2024 (arXiv:2406.11776).
\bibitem{failuremodes} Wynn, et al. Understanding Failure Modes in Multi-Agent Debate. arXiv:2509.05396, 2025.
\bibitem{mmlupro} Wang, Y. et al. MMLU-Pro: A More Robust and Challenging Multi-Task Language Understanding Benchmark. NeurIPS 2024 D\&B.
\bibitem{supergpqa} M-A-P. SuperGPQA: Scaling LLM Evaluation across 285 Graduate Disciplines. 2025.
\bibitem{gsm8k} Cobbe, et al. Training Verifiers to Solve Math Word Problems. arXiv:2110.14168, 2021.
\bibitem{math500} Hendrycks, et al. / Lightman, et al. MATH-500. 2024.
\bibitem{bgem3} Chen, et al. BGE M3-Embedding: Multi-Lingual, Multi-Functionality, Multi-Granularity Text Embeddings. 2024.
\bibitem{robertamnli} Liu, et al. RoBERTa (roberta-large-mnli). 2019.

\end{thebibliography}

\end{document}
